\lecture{2.8}{From DAGs to SWIGs}
\label{lec:2-8}

\subsection{Single World Intervention Graphs}

DAGs represent causal relationships among observed or naturally occurring variables. Potential outcomes such as $\Ya{a}$ do not appear directly on an ordinary DAG. Single World Intervention Graphs, or SWIGs, extend DAGs so that counterfactual variables can be represented graphically.

\begin{definition}[Single World Intervention Graph]
A \emph{Single World Intervention Graph} is a graph obtained from a causal DAG by splitting intervention nodes into a random part and a fixed intervention part, then relabeling downstream variables as counterfactual variables under that intervention.
\end{definition}

"Single world" emphasizes that the graph represents one intervention world at a time. A SWIG for $A=1$ and a SWIG for $A=0$ are distinct intervention worlds.

\subsection{Splitting Treatment Nodes}

Suppose the original DAG is
\[
W\to A\to Y,\qquad W\to Y.
\]
To construct the SWIG for intervention $A=a$, split the node $A$ into:
\begin{itemize}
  \item a random part $A$, representing the treatment the unit would naturally receive;
  \item a fixed part $a$, representing the value imposed by intervention.
\end{itemize}
Incoming arrows into $A$ go to the random part. Outgoing arrows from $A$ leave the fixed part.

\begin{figure}[H]
\centering
\begin{tikzpicture}[node distance=1.8cm,>=Latex]
  \node[draw, circle, fill=blue!5] (W1) {$W$};
  \node[draw, circle, fill=green!5, right=of W1] (A1) {$A$};
  \node[draw, circle, fill=orange!7, right=of A1] (Y1) {$Y$};
  \draw[->, thick] (W1) -- (A1);
  \draw[->, thick] (A1) -- (Y1);
  \draw[->, thick] (W1) to[bend right=25] (Y1);
  \node[right=.7cm of Y1] (arrow) {$\Longrightarrow$};
  \node[draw, circle, fill=blue!5, right=.9cm of arrow] (W2) {$W$};
  \node[draw, circle, fill=green!5, right=of W2] (A2) {$A$};
  \node[draw, rectangle, fill=gray!10, below=.6cm of A2, minimum width=.8cm] (a) {$a$};
  \node[draw, circle, fill=orange!7, right=of A2] (Y2) {$\Ya{a}$};
  \draw[->, thick] (W2) -- (A2);
  \draw[->, thick] (a) -- (Y2);
  \draw[->, thick] (W2) to[bend right=25] (Y2);
\end{tikzpicture}
\caption{A DAG is converted to a SWIG by splitting the intervened treatment node.}
\end{figure}

\subsection{Fixed Intervention Nodes}

The fixed node $a$ is not random. It represents the intervention value. Because it is fixed by the analyst's hypothetical intervention, it has no probability distribution in the intervention world.

The fixed node inherits outgoing arrows from the original treatment node. If $A$ had an arrow to $Y$ in the DAG, then $a$ has an arrow to $\Ya{a}$ in the SWIG. If $A$ had an arrow to a mediator $M$, then $a$ has an arrow to the counterfactual mediator $\M(a)$.

\subsection{Random Counterfactual Nodes}

Downstream nodes are relabeled as counterfactuals. If $Y$ is downstream of $A$, then under intervention $A=a$ the outcome node becomes $\Ya{a}$. If $M$ is downstream of $A$, it becomes $M(a)$.

Nodes not downstream of the intervention keep their original labels. In the graph above, $W$ is a baseline covariate before $A$, so it remains $W$.

\subsection{Constructing SWIGs from DAGs}

The construction rule is:
\begin{enumerate}
  \item Start with a causal DAG.
  \item Choose the node or nodes to intervene on.
  \item Split each intervention node into a random part and fixed part.
  \item Send incoming arrows to the random part.
  \item Send outgoing arrows from the fixed part.
  \item Relabel downstream nodes as counterfactual variables under the intervention.
\end{enumerate}

For a mediation DAG $A\to M\to Y$ with also $A\to Y$, intervening on $A=a$ gives:

\begin{figure}[H]
\centering
\begin{tikzpicture}[node distance=1.7cm,>=Latex]
  \node[draw, circle, fill=green!5] (A) {$A$};
  \node[draw, circle, fill=blue!5, right=of A] (M) {$M$};
  \node[draw, circle, fill=orange!7, right=of M] (Y) {$Y$};
  \draw[->, thick] (A) -- (M);
  \draw[->, thick] (M) -- (Y);
  \draw[->, thick] (A) to[bend right=25] (Y);
  \node[right=.7cm of Y] (arr) {$\Longrightarrow$};
  \node[draw, circle, fill=green!5, right=.9cm of arr] (Ar) {$A$};
  \node[draw, rectangle, fill=gray!10, below=.6cm of Ar, minimum width=.8cm] (a) {$a$};
  \node[draw, circle, fill=blue!5, right=of Ar] (Ma) {$M(a)$};
  \node[draw, circle, fill=orange!7, right=of Ma] (Ya) {$Y(a)$};
  \draw[->, thick] (a) -- (Ma);
  \draw[->, thick] (Ma) -- (Ya);
  \draw[->, thick] (a) to[bend right=25] (Ya);
\end{tikzpicture}
\caption{A mediation DAG and the SWIG under intervention $A=a$.}
\end{figure}

If we intervene on both treatment and mediator, setting $A=a$ and $M=m$, the downstream outcome becomes $Y(a,m)$.

\subsection{Reading Counterfactual Quantities from SWIGs}

SWIGs help identify which counterfactual variables are generated by a particular intervention. In the simple treatment graph, the SWIG for $A=a$ contains $\Ya{a}$, so $\E\{\Ya{a}\}$ is a natural causal parameter.

In the mediation graph, the SWIG for $A=a$ contains $M(a)$ and $Y(a)$. A SWIG for the joint intervention $(A=a,M=m)$ contains $Y(a,m)$. These are different causal questions:
\[
Y(a) \quad\text{allows the mediator to respond naturally to } A=a,
\]
whereas
\[
Y(a,m) \quad\text{sets both treatment and mediator}.
\]

\begin{keyidea}
SWIGs make the intervention explicit. Different interventions on the same original DAG can generate different counterfactual variables and therefore different causal questions.
\end{keyidea}
